\documentclass[10.5pt,a4paper]{article}

\usepackage{fontspec}
\setmainfont{Arial}

\usepackage[french]{babel}
\frenchsetup{StandardItemLabels=true}

\usepackage{geometry}
\geometry{top=2cm, bottom=2cm, left=2.2cm, right=2.2cm}

\usepackage{enumitem}
\setlist[itemize]{leftmargin=1.2em, itemsep=3pt, topsep=2pt, label=\textbullet}

\setlength{\parindent}{0pt}
\setlength{\parskip}{4pt}
\pagestyle{empty}

\begin{document}

\begin{center}
  {\Large\bfseries Synthèse du projet HABS Neuro-AI 2026}\\[4pt]
  {\normalsize Équipe 12 (Mathias Chebbah, Aurélien Saint-Grégoire, Ashley Russo) \textbar\ 23 novembre 2026}
\end{center}

\vspace{8pt}

{\large\bfseries 1. Réponses aux remarques méthodologiques}

\vspace{3pt}

\begin{itemize}
  \item \textbf{Cadrage.} Abandon du storytelling chirurgical. L'étude se concentre sur la neurophysiologie descriptive du stress aigu et sur le monitoring de populations non-verbales (pédiatrie, locked-in, coma).
  \item \textbf{Validité du critère.} Le critère de stress ($V \le 4$, $A \ge 6$) inclut 89\,\% à 97\,\% d'émotions négatives réelles (peur, colère, dégoût). Seuls 5\,\% des essais sont incohérents. L'optimum de détection se situe à $V \le 3$, $A \ge 7$.
  \item \textbf{Spécificité des émotions.} La peur active tout le corps (cœur ralenti, sueur, respiration, chute alpha). Le dégoût ralentit le cœur sans faire transpirer. La tristesse est physiologiquement muette en 10 secondes. Le clustering non supervisé échoue.
  \item \textbf{Cœur et morphologie ECG.} Le cœur freine immédiatement via le nerf vague (+5 ms par battement). La conduction auriculo-ventriculaire reste strictement plate (intervalle PR inchangé), prouvant l'absence d'adrénaline en 10 secondes. Le raccourcissement du QTc était un artefact mathématique.
  \item \textbf{Potentiel évoqué cardiaque (HEP).} L'ICA n'a pas détruit le signal. Le HEP calculé sur données brutes non nettoyées donne exactement le même résultat.
  \item \textbf{Évaluation subjective.} Devant une scène d'horreur, les spectateurs calmes ne montrent aucun ralentissement cardiaque et aucune hausse de sueur. Seule l'attention visuelle est automatique.
\end{itemize}

\vspace{8pt}

{\large\bfseries 2. Ce qu'il faut retenir de l'étude}

\vspace{3pt}

\begin{itemize}
  \item \textbf{L'émotion précède la réaction.} Le corps réagit au ressenti individuel, pas passivement aux images projetées sur l'écran.
  \item \textbf{Un mécanisme purement vagal.} Sur un stimulus court de 10 secondes, le cœur freine par le nerf vague, sans mobilisation sympathique adrénergique.
  \item \textbf{Un trait individuel stable.} La sensibilité au stress est constante d'un jour à l'autre ($r = 0,70$) et se lit déjà au repos par un rythme cardiaque plus rapide.
  \item \textbf{Limite des modèles en 10 secondes.} La détection par IA plafonne à 59\,\% sur de nouveaux sujets (66\,\% face au calme profond), ce qui reste insuffisant pour un usage clinique direct.
\end{itemize}

\vspace{8pt}

{\large\bfseries 3. La cascade neurophysiologique en 10 secondes}

\vspace{3pt}

\begin{itemize}
  \item \textbf{0 à 1 s (Attention réflexe).} Chute brutale des ondes alpha occipitales. Ce réflexe visuel de fixation survient chez tout le monde, y compris chez les personnes impassibles.
  \item \textbf{1 à 3 s (Évaluation cognitive).} Le cerveau analyse la scène. Si la menace n'est pas validée subjectivement, la cascade s'arrête et le corps reste au repos.
  \item \textbf{2 à 5 s (Frein vagal).} Si la peur est confirmée, le nerf vague libère de l'acétylcholine sur le nœud sinusal et ralentit immédiatement le cœur.
  \item \textbf{4 à 7 s (Sudation sympathique).} Activation retardée des glandes sudoripares sous commande sympathique (présente dans la peur, absente dans le dégoût).
  \item \textbf{État basal de repos.} Les personnes sensibles présentent un cœur au repos plus rapide et une activité gamma frontale plus élevée avant même le début du stimulus.
\end{itemize}

\end{document}
